\documentclass[11pt,openany]{book}
\usepackage{robbook}
\SetROBAccent{ROBCyan}
\SetROBVolume{VOLUME 1 - AGES 8-12}{MEET ROB}

\begin{document}
\ROBTitle{Building R.O.B. - Volume 1}{MEET ROB}{A robot is a team of systems}{A picture-book field guide for curious makers ages 8-12}{rob-systems-feedback-lab.png}

\ROBFrontMatter{Welcome, junior systems engineer}{Before we touch a wire, we learn to see the whole robot.}{%
ROB is a real, evolving machine built over many years. These pages use real workshop photographs to explain how a big robot becomes possible: one smaller problem at a time.

This volume is meant to be read with a parent, teacher, or maker mentor. Its activities use paper, pencils, blocks, and small battery-powered parts. The full-size robot contains heavy treads, powerful motors, pinch points, and high-energy batteries. Those parts belong to trained adults with a tested emergency-stop plan.

\begin{center}
\includegraphics[height=0.34\textheight]{rob-systems-feedback-lab.png}\\[-1pt]
{\scriptsize\itshape Original illustrated portrait inspired by the real ROB; use photographs and drawings, not this artwork, as engineering evidence.}
\end{center}
}

\ROBSeriesMap{Read the large words first, study the photographs, then try the mission box. Words in \textbf{bold} belong in your new engineering vocabulary. When a detail is not proven by the archive, the book says so instead of guessing.}
\ROBArduinoStatus
\ROBAmberStatus
\ROBSourceRoot

\begin{safetybox}[title={THE THREE-ZONE RULE}]
\textbf{Green zone:} paper models, unplugged code games, and simulations. Work independently.\par
\textbf{Amber zone:} breadboards, LEDs, sensors, and small low-voltage kits. Work with an adult.\par
\textbf{Red zone:} ROB's batteries, motor wiring, milling, welding, treads, arms, and actuators. Observe from a marked boundary unless a qualified adult has made the machine safe.
\end{safetybox}

\tableofcontents
\clearpage

\mainmatter
\ROBChapterPhoto{Mission 1}{What makes a machine a robot?}{2019-rob-front-corner.jpg}

\chapter{A robot is more than a machine}

Look closely at ROB. You can find parts that roll, parts that bend, parts that see, parts that carry electricity, and parts that run programs. A robot appears to be one creature, but engineers see a \textbf{system of systems}.

A machine becomes robot-like when it can take in information, choose what to do according to rules, and cause a change in the world. That does not mean it is alive or that it thinks like a person. It means its parts cooperate.

\ROBSimpleSystemDiagram

\begin{conceptbox}
The sense-think-act loop is also an \textbf{input-process-output} pattern. A keyboard sends input to a computer. Software processes it. A screen produces output. Feedback tells the next loop what changed.
\end{conceptbox}

\ROBPhoto{2019-rob-side-profile.jpg}{A 2019 side view shows the big idea already in place: a tracked base, a torso, articulated arms, a camera head, wiring, and a computer-controlled personality.}

\begin{missionbox}[title={SYSTEMS HUNT}]
On the photograph, point to one possible \textbf{input}, one \textbf{processor}, and one \textbf{output}. Then name one thing you cannot see but the robot still needs, such as electrical power or software.
\end{missionbox}

\ROBPhoto[0.52\textheight]{2020-rob-front.jpg}{A 2020 front view gives you another systems puzzle. Trace one path from a sensor, through hidden computing, to a part that can move or make sound.}

\clearpage
\section*{Five teams inside ROB}

\begin{tabularx}{\textwidth}{>{\robcondensed\bfseries\color{ROBAccent}}p{1.05in} >{\RaggedRight\arraybackslash}X >{\RaggedRight\arraybackslash}p{2.15in}}
\toprule
TEAM & JOB & EXAMPLES IN ROB\\
\midrule
Structure & Carries loads and keeps parts aligned. & Chassis plates, wooden prototype panels, brackets, shafts.\\
Motion & Changes position by turning or extending. & Treads, flippers, arm joints, torso actuator.\\
Power & Stores and safely distributes energy. & LiFePO4 battery bank described in the build presentation, power buses, converters.\\
Sensing & Measures the robot and its surroundings. & IR distance sensors, IMU, OAK-D depth camera, RPLidar.\\
Control & Turns human intent and sensor data into bounded commands. & Arduino base firmware, Mac mini running Cerebro, phone and wearable controllers.\\
\bottomrule
\end{tabularx}

No team can succeed alone. A powerful motor without a bracket may twist itself loose. A perfect camera without software only produces pixels. A clever program without an emergency stop can be dangerous.

\begin{buildbox}[title={MAKE A SYSTEMS MOBILE}]
Cut five paper circles and label them Structure, Motion, Power, Sensing, and Control. Tie them to one ring. Add a sixth card labeled Human. Tug one card and watch the whole mobile move: a change in one subsystem affects the others.
\end{buildbox}

\ROBChapterPhoto{Mission 2}{Structure carries every idea}{2021-shoulder-gear-and-cable.jpg}

\chapter{The skeleton and the shell}

The \textbf{chassis} is the robot's load-carrying frame. ROB's archive shows wooden prototype panels, metal brackets, shafts, bearings, fasteners, tread frames, and later mechanical revisions. Prototypes are useful because they let a builder test size and motion before paying for final parts.

ROB began with one hard mechanical promise: the purchased SuperDroid Robots \textbf{LT2 tread system} had to fit. Instead of immediately cutting expensive material, the builder copied the tread envelope in cardboard. Cardboard made it cheap to move a wall, change a hole, test a battery-sized box, and ask whether the flipper could sweep without striking the body. The useful shapes could then become templates for maple plywood. A playful three-wheel idea, inspired by the famous movie robot Johnny 5, explored whether a sideways wheel could reduce tread scrub during a turn.

\begin{fieldnotebox}[title={CARDBOARD IS A QUESTION-MAKING TOOL}]
A prototype does not need to be strong enough for the final robot. It needs to answer the question written on it. Mark every model with its question, scale, date, and result: \emph{Does the LT2 fit? Can the battery be removed? Does the flipper clear? Can a hand reach the connector?} A failed cardboard shape can save a sheet of plywood, a machined plate, and a week of work.
\end{fieldnotebox}

The tracked base must do several jobs at once:

\begin{itemize}
  \item hold the robot's weight without bending too much;
  \item keep sprockets, wheels, and chains aligned;
  \item protect batteries and wiring;
  \item provide openings for inspection and repair;
  \item leave room for a central actuator and moving flippers.
\end{itemize}

\ROBPhotoTwo{2021-cable-loop-and-arm.jpg}{A cable loop beside an arm shows that structure must provide a safe path for things that bend and move.}{2021-shoulder-gear-and-cable.jpg}{A shoulder gear, plate, and cable share one crowded interface. Structure is also about clearance and service access.}

\begin{conceptbox}
A measurement is data. A useful measurement has a \textbf{value}, a \textbf{unit}, and an idea of \textbf{uncertainty}. Writing ``48'' is incomplete. Writing ``48.0 mm, measured with a caliper'' can guide another maker.
\end{conceptbox}

\begin{buildbox}[title={PAPER CHASSIS TEST}]
Build two bridges from one sheet of paper. Leave one flat. Fold the other into a channel or triangle. Span the same gap and add coins one at a time. Record how many coins each design holds. Shape can add stiffness without adding material.
\end{buildbox}

\begin{buildbox}[title={MAKE A FULL-SIZE INTERFACE MOCKUP}]
Choose a small robot wheel or print a paper circle at full size. Tape it to a cardboard side panel, then add paper envelopes for the motor, battery, and cable bend. Move the parts before cutting any real material. Photograph the final mockup beside a ruler and write three things it proved and three things it could not prove.
\end{buildbox}

\begin{fieldnotebox}[title={ROB IS A NAME, NOT A DIMENSION SET}]
ROB is a friendly name taken from the word ``robot''; it is not an acronym. Rodolfo Aramayo began ROB after leaving systems-software research work at Verizon Labs and built the early body by hand in College Station, Texas. He dates the new treaded prototype to approximately 2016, before the later metal revisions. The present height, width, length, ready-to-run mass, and center of mass have not been measured for publication. Never scale a photograph to copy the machine: a future reproduction begins with its own measured drawing, stability study, and adult engineering review.
\end{fieldnotebox}

\ROBPhoto[0.53\textheight]{2020-rob-overhead.jpg}{The overhead view turns ROB into a structure puzzle: find the plates, mounts, gaps, and cable paths that let many systems share one body.}

\clearpage
\section*{Iteration is not failure}

\ROBPhotoTwo{2021-open-electronics-bay.jpg}{An open body reveals how structure must support electronics while leaving room to inspect and repair them.}{2021-shoulder-gear-and-cable.jpg}{Parts that rotate need clearance. Cables need bend room. A successful structure makes both paths understandable.}

An early design can answer a question even when it is not the final design. The ROB presentation records revisions such as moving the flippers, changing brackets, adding a third-wheel idea for easier turning, and finding new spaces for sensors and computing.

Engineers often work in this loop:

\begin{center}
\begin{tikzpicture}[node distance=8mm,every node/.style={font=\robcondensed\bfseries\small,align=center,draw=ROBSteel,rounded corners=1mm,minimum width=1.07in,minimum height=0.48in}]
\node[fill=ROBCyan!9] (ask) {ASK};
\node[right=of ask,fill=ROBAmber!10] (sketch) {SKETCH};
\node[right=of sketch,fill=ROBGreen!9] (build) {BUILD};
\node[right=of build,fill=ROBPurple!8] (test) {TEST};
\node[right=of test,fill=ROBRed!7] (learn) {LEARN};
\draw[->,very thick,ROBAccent] (ask)--(sketch);
\draw[->,very thick,ROBAccent] (sketch)--(build);
\draw[->,very thick,ROBAccent] (build)--(test);
\draw[->,very thick,ROBAccent] (test)--(learn);
\draw[->,very thick,ROBAccent,bend left=32] (learn.south) to (ask.south);
\end{tikzpicture}
\end{center}

The last arrow is the secret: learning returns you to a better question.

\ROBChapterPhoto{Mission 3}{Power and information are different}{rob-systems-feedback-lab.png}

\chapter{Energy makes motion possible}

Electric circuits carry energy and information, but those are not always the same job. A thick red cable may carry enough current to drive a motor. A small signal wire may only tell the motor controller what speed to use.

ROB's presentation describes four 20 Ah LiFePO4 batteries arranged across the base for a total stated capacity of 80 Ah. A 2024 photograph shows two large batteries with red power leads inside the chassis. The archive does not prove the final series/parallel arrangement, protection devices, or present-day capacity, so this book does not invent them.

The builder's newest record describes the negative leads meeting at a return bus and the positive leads passing through the front power switch before reaching a 12 V distribution bus. Other converters create 24 V and 48 V domains for different jobs. That story is a map for an adult to verify, not a circuit for a young reader to copy. Amp-hours describe charge capacity; they do not tell us the voltage, safe current, or runtime by themselves.

\begin{safetybox}
Large battery banks can produce fire, burns, arc flash, toxic smoke, and violent motion. Never copy a photograph as a wiring diagram. A qualified adult must design fusing, a disconnect, charging, cable sizes, insulation, strain relief, and the physical emergency-stop circuit.
\end{safetybox}

\ROBPhoto[0.46\textheight]{rob-systems-feedback-lab.png}{Conceptual illustration: amber paths represent energy and cyan paths represent information. It is a thinking map, not ROB's as-built wiring diagram.}

\begin{missionbox}[title={ENERGY OR MESSAGE?}]
Sort these into likely energy paths or information paths: battery cable, USB cable, motor shaft, camera image, PWM control wire, speaker wire, emergency-stop contact. Some can carry both. Explain your choice instead of guessing by color. In \emph{Circuits and Signals}, a supervised Maker Faire lab opens a sacrificial, unplugged USB cable and uses its verified 5 V and GND conductors to power a small toy. It turns this sorting question into a real circuit while keeping cutting and electrical verification in an adult-only station.
\end{missionbox}

\ROBChapterPhoto{Mission 4}{Sensors turn the world into numbers}{2025-rob-face-portrait.jpg}

\chapter{How ROB notices things}

A sensor is a translator. It converts something physical into data a computer can use.

\begin{itemize}
  \item The Arduino file declares six Sharp IR sensor objects. The active warning routine samples only the front and rear pairs; the left/right pair appears only in commented telemetry.
  \item An MPU9250 \textbf{IMU} measures acceleration, turning rate, and magnetic direction. Software combines those readings into yaw, pitch, and roll.
  \item An OAK-D camera can provide a color image plus depth values in millimeters.
  \item RPLidar reports laser samples around the robot so software can reason about obstacles and maps.
\end{itemize}

\ROBPhoto{2025-rob-face-portrait.jpg}{ROB's speaker face, central depth camera, and upper mechanisms show that different sensors answer different questions; no single sensor knows everything.}

\begin{fieldnotebox}
The build presentation records a real lesson: IR light sometimes bounced from the floor, and grass produced noisy obstacle reports. In the checked-in Arduino code, sensor reading and warnings still run, but the lines that set the forward/backward shutdown flags are commented out. That is an experimental condition, not a finished safety system.
\end{fieldnotebox}

\begin{buildbox}[title={NOISY SENSOR GAME}]
One student secretly chooses a distance from 10 to 100 cm. A ``sensor'' student reports it five times, adding or subtracting a few centimeters each time. Should a robot stop after one low reading, after several low readings, or after combining another sensor? Discuss the tradeoff between false alarms and missed obstacles.
\end{buildbox}

\ROBPhoto[0.50\textheight]{2020-rob-front.jpg}{Sensors ride on a moving, changing machine. A useful reading depends on where the sensor is mounted, which way it points, what the robot is doing, and what surrounds it.}

\ROBChapterPhoto{Mission 5}{The brain is a conversation}{rob-systems-feedback-lab.png}

\chapter{Commands travel in layers}
\ROBSourceNow{Cerebro/Cerebro/ROBSerialBox.m}{This implementation connects the Mac coordinator to the present Base Arduino serial protocol; young readers can follow it as evidence that one visible action crosses several systems.}

ROB's Mac mini runs the Cerebro app. Controllers on a phone, Apple Watch, or Vision Pro can express human intent. Cerebro checks who is allowed to control, keeps track of a live session, coordinates cameras and higher-level behavior, and sends a compact command toward the base Arduino.

\ROBControlStackDiagram

Each layer has a smaller job:

\begin{enumerate}
  \item The human chooses a goal, such as ``drive slowly forward.''
  \item A controller measures the joystick and creates numbers.
  \item The network proves which paired device sent the message.
  \item Cerebro decides who currently has control and translates the request.
  \item USB serial carries a fixed-width text frame to the Arduino.
  \item The Arduino sets brake, direction, PWM, and actuator commands.
\end{enumerate}

\begin{conceptbox}
This is \textbf{abstraction}: each layer offers a simpler promise to the layer above. The phone does not need to know which Arduino pin drives the left tread. The Arduino does not need to know whether the command began on a phone or a headset.
\end{conceptbox}

\clearpage
\section*{What happens when a message disappears?}

For a moving robot, silence must not mean ``keep going forever.'' ROB's layers use different time rules: Vision uses a short input lease, the Watch relay has its own watchdog, Cerebro requires a fresh controller snapshot, and the Arduino has a separate heartbeat that commands zero speed if updates stop. Their physical stopping behavior still has to be measured.

\ROBStopChainDiagram

\begin{safetybox}
Software stop layers are useful, but they do not replace a physical emergency stop that removes motion authority independently. Sensors can fail. Programs can freeze. Networks can disappear. Safety works best in layers with different failure modes.
\end{safetybox}

\begin{missionbox}[title={MESSAGE RELAY}]
Form a line of five people: Operator, Controller, Cerebro, Arduino, Motor. Pass a card marked \texttt{FORWARD 20}. Every person must receive a fresh card before a five-count expires. If a card is delayed or unreadable, the next person passes \texttt{STOP}. This is a human-sized watchdog demonstration.
\end{missionbox}

\ROBChapterPhoto{Mission 6}{A robot grows through versions}{2025-rob-side-full-body.jpg}

\chapter{ROB's build timeline}
\ROBSourceNow{Presentation/ROB_v3.pdf}{This historical design deck and the dated gallery photographs are the evidence for the timeline. It is documentation evidence, not executable source code.}

The archive is not a perfect diary, but it shows a machine changing in public.

\section*{2019: the whole idea moves}
\ROBPhotoTwo{2019-rob-front-corner.jpg}{The early black tracked robot already combines camera eyes, arms, a body, and a mobile base.}{2019-rob-side-profile.jpg}{The side view makes weight, balance, tread geometry, and exposed mechanisms easier to study.}

\section*{2020-2021: open the body and learn}
\ROBPhotoTwo{2020-rob-overhead.jpg}{An overhead view reveals how many mechanical and electrical paths compete for space.}{2021-open-electronics-bay.jpg}{An open rear bay exposes loose wiring beside an arm. Access for service is part of the design.}

\clearpage
\section*{2022: add computers, cameras, and brackets}
\ROBPhoto{2022-chassis-wiring-overhead.jpg}{This 2022 overhead build view records the year's integration work: motor hardware, wiring, mounting pieces, and tread chassis all compete for space. This photograph is intentionally shared with the circuits book: Volume 1 treats it as timeline evidence, while Volume 2 studies its routing.}

\section*{2024: measure and rebuild the drivetrain}
\ROBPhoto{2024-nearly-assembled-base.jpg}{The nearly assembled 2024 base records a major physical revision. Volume 3 studies its mechanics in detail; here it is one dated step in ROB's longer story.}

\clearpage
\section*{2025: refine the current character}
\ROBPhotoTwo{2025-rob-side-full-body.jpg}{The current body, tracked base, camera head, and articulated arm show how ROB's identity grew through many revisions.}{2025-neck-gear-train.jpg}{The neck mechanism exposes stacked gears, supports, and a maker's willingness to keep iterating.}

\begin{fieldnotebox}[title={A FAIR PRESENTATION IS A TEST}]
Maker Faire adds new requirements: a safe operating boundary, reliable setup, clear explanations, charging and transport plans, consent around cameras, and a quick way to stop motion. A robot is not finished when it moves; it is ready when people can experience it safely.
\end{fieldnotebox}

\chapter{Your first engineering log}

Choose an imaginary helper robot. Give it one job, such as carrying library books or watering a plant. Fill in this one-page design brief.

\begin{tabularx}{\textwidth}{>{\robcondensed\bfseries\color{ROBAccent}}p{1.45in} X}
\toprule
USER & Who needs help? \rule{0pt}{0.42in}\\
TASK & What one job should the robot do? \rule{0pt}{0.42in}\\
INPUTS & What must it sense or be told? \rule{0pt}{0.42in}\\
OUTPUTS & What motion, light, or sound will it make? \rule{0pt}{0.42in}\\
POWER & What safe classroom energy source will it use? \rule{0pt}{0.42in}\\
STOP RULE & What event must immediately make it stop? \rule{0pt}{0.42in}\\
TEST & How will you know the idea worked? \rule{0pt}{0.42in}\\
\bottomrule
\end{tabularx}

\begin{buildbox}[title={DRAW THE BLOCK DIAGRAM}]
Draw boxes for Sense, Think, and Act. Add arrows showing the information path. Then add a Human box and a Stop path. A clear diagram is a tool for thinking, not decoration.
\end{buildbox}

\ROBChapterPhoto{Deep Lab}{See the loop, not just the parts}{rob-systems-feedback-lab.png}

\chapter{Think like a systems detective}

A list of parts tells us what a robot \emph{has}. A system model tells us what the parts must accomplish together. Systems engineers begin with a boundary: what is inside the system, what is outside, and what may cross the boundary? Energy, information, material, force, heat, and sound can cross it. Drawing the boundary prevents a common mistake---explaining one interesting part while forgetting the whole job.

\begin{missionbox}[title={THE WATER-DELIVERY ROBOT}]
Imagine a small classroom robot whose only job is to carry one sealed cup from a filling station to a table. Name the user and success condition. Then list what crosses its boundary: electrical energy enters from a battery; destination information enters from a person; wheel forces act on the floor; heat and sound leave; the cup enters and exits. Which crossing could make the mission fail?
\end{missionbox}

\section*{A feedback loop has memory}

Open-loop action says, ``Run the motor for one second.'' Closed-loop action says, ``Measure the result, compare it with the goal, and choose the next action.'' The second approach needs a remembered goal and fresh measurements. It can correct some disturbances, such as a wheel moving less on carpet, but only when the sensor measures something useful.

\begin{center}
\begin{tikzpicture}[node distance=8mm,>=stealth,every node/.style={font=\robcondensed\small,align=center}]
\node[draw=ROBAccent,fill=ROBAccent!8,rounded corners,minimum width=1.15in,minimum height=.48in] (goal) {GOAL};
\node[draw=ROBPurple,fill=ROBPurple!8,rounded corners,right=of goal,minimum width=1.15in,minimum height=.48in] (compare) {COMPARE};
\node[draw=ROBGreen,fill=ROBGreen!8,rounded corners,right=of compare,minimum width=1.15in,minimum height=.48in] (act) {ACT};
\node[draw=ROBCyan,fill=ROBCyan!8,rounded corners,below=9mm of compare,minimum width=1.15in,minimum height=.48in] (sense) {SENSE};
\draw[->,thick] (goal)--(compare); \draw[->,thick] (compare)--(act);
\draw[->,thick] (act.east) -- ++(.45,0) |- (sense.east);
\draw[->,thick] (sense)--(compare);
\end{tikzpicture}
\end{center}

\begin{conceptbox}[title={A SENSOR DOES NOT KNOW}]
A sensor translates a physical effect into a signal. A camera produces pixels; it does not automatically know ``cup.'' Software builds an estimate from measurements, assumptions, and prior information. Every estimate has limits. Strong glare, a hidden object, a dirty lens, or an old timestamp can change what the data means.
\end{conceptbox}

\section*{Make three predictions}

For every model, predict a normal case, a changed case, and a failure case. If ROB approaches a cardboard box, what should a distance reading do? What changes if the box is dark or angled? What should happen if readings stop arriving? A prediction turns ``I think it works'' into something observable.

\begin{tabularx}{\textwidth}{>{\robcondensed\bfseries\color{ROBAccent}}p{1.2in} X X}
\toprule
CASE & PREDICTION & EVIDENCE TO COLLECT\\
\midrule
Normal & What pattern should appear? & A number, timestamp, photo, or state.\\
Changed & Which variable changes while others stay fixed? & Before-and-after observations.\\
Failure & Which safe state should appear? & Timeout, rejection, and recovery record.\\
\bottomrule
\end{tabularx}

\clearpage
\chapter{Uncertainty is useful information}

Measurements are not perfect labels. A ruler has limited marks. A camera pixel covers an area. A distance sensor samples a cone or ray. Repeating a measurement reveals spread. Engineers report a value with its unit, method, revision, and uncertainty instead of hiding variation.

\begin{buildbox}[title={PAPER SENSOR MAP}]
Place a sheet of graph paper on a table and choose one square as a pretend robot. A partner secretly places three paper obstacles. You may ask for eight straight-line distance measurements in chosen directions. Mark each result, including ``no return.'' Draw two maps: space supported by evidence and space still unknown. Compare your map with the real layout. No electronics are needed.
\end{buildbox}

\begin{advancedbox}[title={CAUSE, CORRELATION, AND A FAIR TEST}]
If a robot slows down when its battery reading falls, the two events are correlated. That alone does not prove the battery caused the slowdown; floor friction, software limits, temperature, or load may also have changed. A fair test changes one planned variable, controls the others as well as practical, repeats trials, and records exceptions.
\end{advancedbox}

\begin{safetybox}
These reasoning tools apply to ROB, but youth investigations remain on paper, in simulation, or with current-limited classroom kits. Never probe ROB's batteries, power branches, motor wiring, moving treads, actuators, or exposed mechanisms.
\end{safetybox}



\clearpage
\section*{Field words}

\begin{multicols}{2}
\ROBFact{Actuator}{A device that causes motion.}\par
\ROBFact{Chassis}{The load-carrying frame.}\par
\ROBFact{Command}{A message asking a system to do something.}\par
\ROBFact{Feedback}{A measurement of what happened after an action.}\par
\ROBFact{Input}{Information entering a system.}\par
\ROBFact{Iteration}{A new version made after learning from a test.}\par
\ROBFact{Output}{Information or action produced by a system.}\par
\ROBFact{Prototype}{A version built to answer questions.}\par
\ROBFact{Sensor}{A device that translates a physical quantity into data.}\par
\ROBFact{System}{Connected parts that cooperate for a purpose.}\par
\ROBFact{Watchdog}{A timer that responds when expected updates stop.}\par
\end{multicols}

\begin{missionbox}[title={FINAL CHECK}]
Can you explain why ROB needs both energy paths and information paths? Can you name three systems inside him? Can you explain why a watchdog and a physical emergency stop are different? If yes, you are ready for Volume 2.
\end{missionbox}

\vfill
\ROBSource{Presentation/ROB\_v3.pdf; ORobotics private gallery originals and gallery metadata; ROBArduino Base, retired Torso, and retired Head sketches; current Cerebro, ROBController, and ROBControllerVision source. Arduino archive inspected 2026-08-10.}

\end{document}
